% Part: first-order-logic
% Chapter: model-theory
% Section: nonstandard-arithmetic

\documentclass[../../../include/open-logic-section]{subfiles}

\begin{document}

\olfileid{mod}{bas}{nsa}
\section{পাটিগণিতের অমানক মডেল}

\begin{defn}
পাটিগণিতের !!{language}~$\Lang{L_N}$-এ আছে একটি !!{constant}~$\Obj{0}$,
একটি দ্বিস্থানী !!{predicate}~$<$, একটি একস্থানী !!{function}~$\prime$,
এবং দুটি দ্বিস্থানী !!{function}s~$+$ ও~$\times$।
\begin{enumerate}
\item পাটিগণিতের \emph{মানক মডেল} হলো $\Lang{L_N}$-এর !!{structure}
  $\Struct{N}$; এতে $\Nat = \{ 0, 1, 2, \dots\}$, এবং
  $\Obj{0}$-এর ব্যাখ্যা $0$, $<$-এর ব্যাখ্যা $\Nat$-এর উপর ক্ষুদ্রতর
  সম্পর্ক, আর $\prime$, $+$ ও~$\times$-এর ব্যাখ্যা যথাক্রমে $\Nat$-এর
  উপর উত্তরসূরি, যোগ ও গুণ।
\item \emph{সত্য পাটিগণিত} হলো তত্ত্ব~$\Theory{N}$।
\end{enumerate}
\end{defn}

$\Lang{L_N}$-এ কাজ করার সময় $\Obj{0}$-এর উপর উত্তরসূরি অপেক্ষকটি
$n$ বার প্রয়োগে পাওয়া $\Obj{0}^{\prime\dots\prime}$ রূপের প্রতিটি
পদকে সংক্ষেপে~$\num{n}$ লিখি।

\begin{defn}
$\Lang{L_N}$-এর কোনো !!{structure}~$\Struct{M}$ \emph{মানক} যদি এবং
কেবল যদি $\Struct{N} \iso \Struct{M}$।
\end{defn}

\begin{thm}\ollabel{thm:non-std}
সত্য পাটিগণিতের !!{enumerable} অমানক মডেল আছে।
\end{thm}

\begin{proof}
$\Lang{L_N}$-এ একটি নতুন !!{constant}~$c$ যোগ করে নিচের তত্ত্বটি
বিবেচনা করো:
\[
\Theory{N} \cup \Setabs{\num{n} < c}{n \in \Nat}.
\]
তত্ত্বটির প্রতিটি সসীম অংশের মডেল আছে। তাই সংহতি উপপাদ্য অনুসারে
এর একটি মডেল~$\Struct{M}$ আছে; নিম্নগামী লোয়েনহাইম--স্কোলেম
উপপাদ্য অনুসারে মডেলটি !!{enumerable} নেওয়া যায়। $\Domain{M}$ যদি
$\Struct{M}$-এর সংজ্ঞাক্ষেত্র হয়, তবে $\Struct{M}$-এ
$\Lang{L_N}$-এর অযৌক্তিক
চিহ্নগুলির ব্যাখ্যা লিখি $\mathbf{z} = \Assign{\Obj{0}}{M} \in
\Domain{M}$, ${\prec} = \Assign{<}{M} \subseteq \Domain{M}^2$,
$* = \Assign{\prime}{M} \colon \Domain{M} \to \Domain{M}$,
এবং $\oplus = \Assign{+}{M}, \otimes = \Assign{\times}{M}
\colon \Domain{M}^2 \to \Domain{M}$। প্রতিটি $x \in \Domain{M}$-এর
জন্য $x$-এর উপর $*$ প্রয়োগে পাওয়া $\Domain{M}$-এর উপাদানটিকে $x^*$ লিখি।
% BN-SRC-507 TARGET-CORRECTION-BEGIN
\par\smallskip\noindent\textnormal{\small\textbf{উৎস-সংশোধন (BN-SRC-507):} হিমায়িত প্রমাণে নির্দিষ্ট পাটিগাণিতিক ভাষা ঘোষণার পরেও একবার সাধারণ ভাষা লেখা হয়েছে এবং সম্পর্কটির সংজ্ঞাক্ষেত্রের বর্গের বদলে গঠনের নামের বর্গ লেখা হয়েছে। এখানে ঘোষিত ভাষা ও সংজ্ঞাক্ষেত্রের সংকেত পুনরুদ্ধার করা হয়েছে।}
% BN-SRC-507 TARGET-CORRECTION-END
% BN-SRC-508 TARGET-CORRECTION-BEGIN
\par\smallskip\noindent\textnormal{\small\textbf{উৎস-সংশোধন (BN-SRC-508):} হিমায়িত গদ্য ব্যাখ্যায় বিধেয় ও অপেক্ষকসহ সব অযৌক্তিক চিহ্নকে “ধ্রুবক” বলা হয়েছে। বাংলা বাক্যে তাদের যথাযথ সমষ্টিগত নাম “চিহ্ন” ব্যবহৃত হয়েছে।}
% BN-SRC-508 TARGET-CORRECTION-END

এবার $h$ যদি $\Struct{N}$ ও~$\Struct{M}$-এর সমরূপতা হতো, তবে
কোনো $n \in \Nat$-এর জন্য $h(n) = \Assign{c}{M}$ হতো। তাই
$\Struct{N}$-এ এমন যেকোনো আরোপ~$s$ নিই, যাতে $s(x)=n$। তখন
$\Sat{N}{\eq[\num{n}][x]}[s]$; \olref[iso]{thm:isom}-এর প্রমাণ অনুসারে
$\Sat{M}{\eq[\num{n}][x]}[h \circ s]$-ও। ফলে $\Assign{c}{M} =
\mathbf{z}^{*\cdots *}$, যেখানে $*$ ঠিক $n$ বার প্রয়োগ করা হয়েছে।
কিন্তু এটি অসম্ভব: অনুমান অনুসারে $\Sat{M}{\num{n} < c}$, আর
$\prec$ অপরাবর্তী। সুতরাং $\Struct{M}$ অমানক।
\end{proof}

\begin{prob}
সেট~$X$-এর উপর সম্পর্ক~$R$ \emph{সুভিত্তিক} যদি এবং কেবল যদি
$R$-এ কোনো অসীম অবরোহী শৃঙ্খল না থাকে; অর্থাৎ $X$-এ এমন
$x_0$, $x_1$, $x_2$,~\dots নেই যাতে $\dots x_2Rx_1Rx_0$।
জার্মেলো--ফ্রেঙ্কেল সেটতত্ত্ব~$ZF$ সঙ্গত ধরে দেখাও, $ZF$-এর
অসুভিত্তিক মডেল আছে; অর্থাৎ এমন মডেল~$\mathfrak{M}$ আছে, যাতে
$\dots x_2 \in x_1 \in x_0$।
\end{prob}

\olref{thm:non-std}-এর অমানক মডেল~$\Struct{M}$ মানক মডেলটির
সঙ্গে প্রাথমিক সমতুল্য। তাই $\Struct{M}$-এর বেশ কিছু ধর্ম
নির্ণয় করা যায়। এই অংশের বাকি আলোচনা সেই কাজেই নিবেদিত;
এর মাধ্যমে $\Theory{N}$-এর !!{enumerable} অমানক মডেলগুলিকে নির্ভুলভাবে
চিহ্নিত করতে পারব।

\begin{enumerate}
\item $\Domain{M}$-এর কোনো উপাদান নিজেই নিজের চেয়ে $\prec$-ছোট নয়:
  $\lforall[x][\lnot x < x]$ বাক্যটি $\Struct{N}$-এ সত্য,
  তাই $\Struct{M}$-এও সত্য।
\item একই ধরনের যুক্তিতে পাই, $\prec$ হলো $\Domain{M}$-এর উপর
  একটি \emph{রৈখিক ক্রম}; অর্থাৎ একটি সম্পূর্ণ, অপরাবর্তী ও
  সংক্রমণশীল সম্পর্ক; এই তিনটি ধর্মই $\Domain{M}$-এর উপর প্রযোজ্য।
\item $\Domain{M}$-এ $\prec$-ক্রমে ক্ষুদ্রতম উপাদানটি হলো
  $\mathbf{z}$।
\item $\Domain{M}$-এর যেকোনো উপাদান তার $*$-উত্তরসূরির চেয়ে
  $\prec$-ছোট; আর $x$-এর চেয়ে বড় উপাদানগুলির মধ্যে
  $\Domain{M}$-এ $\prec$-ক্ষুদ্রতম হলো $x^*$।
\item $\Struct{M}$-এ $\Nat$-এর সমরূপ একটি $\prec$-প্রারম্ভিক
  অংশ আছে: $\mathbf{z}, \mathbf{z}^*, \mathbf{z}^{**}, \dots$।
  একে $\Domain{M}$-এর \emph{মানক অংশ} বলি। $\Domain{M}$-এর
  অন্য যেকোনো উপাদান \emph{অমানক}। $\Domain{M}$-এ
  অমানক উপাদান থাকতেই হবে;
  নইলে \olref{thm:non-std}-এর প্রমাণের অপেক্ষক~$h$ একটি
  সমরূপতা হতো। $n, m, \dots$-কে $\Struct{M}$-এর এই মানক
  অংশে বিচরণশীল !!{variable}s হিসেবে ব্যবহার করি।
\item প্রত্যেক অমানক উপাদান প্রত্যেক মানক উপাদানের চেয়ে বড়।
  কারণ প্রতিটি $n \in \Nat$-এর জন্য
  \[
  \Sat{N}{\lforall[z][(\lnot(\eq[z][\Obj{0}] \lor
  \dots \lor \eq[z][\num{n}]) \lif \num{n} < z)]},
  \]
  তাই $z \in \Domain{M}$ যদি সব মানক উপাদান থেকে ভিন্ন হয়,
  তবে তাদের প্রত্যেকটির চেয়ে তাকে \emph{বড়} হতে হবে।
\item $\mathbf{z}$ ছাড়া $\Domain{M}$-এর প্রতিটি উপাদান,
  $\Domain{M}$-এর কোনো অনন্য উপাদানের $*$-উত্তরসূরি;
  তার পূর্বসূরিকে $^*x$ লিখি। $x=y^*$ হলে $x$ বা~$y$-এর
  একটি মানক হলে দুটিই মানক; একটি অমানক হলে দুটিই অমানক।
\item $\Domain{M}$-এর উপর $\approx$ সমতুল্যতা-সম্পর্কটি এভাবে
  সংজ্ঞায়িত করি: কোনো \emph{মানক}~$n$-এর জন্য $x \oplus n = y$
  অথবা $y \oplus n =x$ হলে এবং কেবল তখনই $x \approx y$।
  অন্যভাবে, $x$ ও~$y$-এর দূরত্ব সসীম হলে এবং কেবল তখনই
  $x \approx y$। $n$ ও~$m$ মানক হলে $n \approx m$। $x$-এর
  \emph{খণ্ড} হলো সমতুল্যতা-শ্রেণি $[x] = \Setabs{y}{x
  \approx y}$।
\item ধরা যাক, $x \prec y$ কিন্তু $x \not\approx y$। যেহেতু
  $\Sat{N}{\lforall[x][\lforall[y][(x < y \lif (x' < y \lor x' =
      y))]]}$, তাই হয় $x^* \prec y$, নয় $x^*=y$। দ্বিতীয়টি
  অসম্ভব, কারণ তাতে $x\approx y$ হয়; সুতরাং $x^* \prec y$।
% BN-SRC-509 TARGET-CORRECTION-BEGIN
  \par\smallskip\noindent\textnormal{\small\textbf{উৎস-সংশোধন (BN-SRC-509):} দুটি ক্ষেত্রের দ্বিতীয়টি বাদ দেওয়ার পর হিমায়িত প্রমাণটি আগেই অনুমিত অসমতাটি পুনরাবৃত্তি করেছে। এখানে প্রথম ক্ষেত্রের শক্তিশালী সিদ্ধান্ত পুনরুদ্ধার করা হয়েছে।}
% BN-SRC-509 TARGET-CORRECTION-END
  একইভাবে $x \prec y$ ও~$x \not\approx y$ হলে
  $x \prec {^*y}$। অতএব $x \prec y$ ও~$x \not\approx y$ হলে
  প্রত্যেক $w \approx x$ প্রত্যেক $v \approx y$-এর চেয়ে
  $\prec$-ছোট। সেই অনুযায়ী, প্রতিটি অমানক খণ্ড~$[x]$ একটি
  উভয় দিকে অসীম শৃঙ্খল গঠন করে:
% BN-SRC-510 TARGET-CORRECTION-BEGIN
  \par\smallskip\noindent\textnormal{\small\textbf{উৎস-সংশোধন (BN-SRC-510):} হিমায়িত বাক্যে প্রত্যেক খণ্ডকে উভয় দিকে অসীম বলা হয়েছে। মানক খণ্ডে ক্ষুদ্রতম উপাদান আছে এবং তার ক্রমপ্রকার স্বাভাবিক সংখ্যার মতো; এই দাবি কেবল অমানক খণ্ডে সত্য।}
% BN-SRC-510 TARGET-CORRECTION-END
  \[
  \cdots \prec  {^{**}x} \prec {^*}x \prec x \prec x^* \prec x^{**}
  \prec \cdots
  \]
  একে $Z$-শৃঙ্খল বলা হয়, কারণ এর ক্রমপ্রকার পূর্ণসংখ্যার ক্রমপ্রকার।
\item $\prec$-ক্রমটি খণ্ডগুলির উপরেও আরোপ করা যায়: $x \prec y$
  এবং তারা ভিন্ন খণ্ডে থাকলে $x$-এর খণ্ড $y$-এর খণ্ডের চেয়ে ছোট।
% BN-SRC-511 TARGET-CORRECTION-BEGIN
  \par\smallskip\noindent\textnormal{\small\textbf{উৎস-সংশোধন (BN-SRC-511):} মূল বাক্যে শুধু প্রতিনিধি দুটি ক্রমে থাকলেই তাদের খণ্ডগুলির মধ্যে কঠোর ক্রম দাবি করা হয়েছে। একই খণ্ডের দুটি ভিন্ন প্রতিনিধিও ক্রমে থাকতে পারে; তাই ভিন্ন খণ্ডের শর্ত স্পষ্ট করা হয়েছে।}
% BN-SRC-511 TARGET-CORRECTION-END
  কোনো খণ্ডে অমানক উপাদান থাকলে সেটি \emph{অমানক} খণ্ড।
  মানক খণ্ডটি ক্ষুদ্রতম খণ্ড।
\item কোনো ক্ষুদ্রতম অমানক খণ্ড নেই: $y$ অমানক হলে এমন
  $x \prec y$ আছে, যাতে $x$-ও অমানক এবং $x \not\approx y$।
  প্রমাণ: মানক মডেল~$\Struct{N}$-এ প্রতিটি সংখ্যাকে দুই দিয়ে
  ভাগ করা যায়, সম্ভবত এক ভাগশেষ রেখে; অর্থাৎ
  $\Sat{N}{\lforall[y][\lexists[x][(y = x + x \lor y = x + x +
        \Obj{0}')]]}$।
% BN-SRC-512 TARGET-CORRECTION-BEGIN
  \par\smallskip\noindent\textnormal{\small\textbf{উৎস-সংশোধন (BN-SRC-512):} হিমায়িত ভাগ-সূত্রে অর্ধেকের চলককে সার্বিক পরিমাণসূচকে বাঁধায় বাক্যটি মিথ্যা। প্রতিটি সংখ্যার জন্য কোনো একটি অর্ধেক থাকার অস্তিত্বসূচক পুনরুদ্ধার করা হয়েছে।}
% BN-SRC-512 TARGET-CORRECTION-END
  প্রাথমিক সমতুল্যতা অনুসারে, প্রতিটি $y \in \Domain{M}$-এর জন্য
  এমন $x \in \Domain{M}$ আছে যাতে হয় $x \oplus x = y$,
  নয় $x \oplus x \oplus \mathbf{z}^*= y$। $x$ মানক হলে
  $y$-ও মানক হতো; তাই $x$ অমানক। উপরন্তু $x$ ও~$y$
  ভিন্ন খণ্ডে আছে; অর্থাৎ $x \not\approx y$। এটি দেখতে
  বিপরীত ধরে নিই: কোনো মানক~$n$-এর জন্য $x \oplus n = y$।
  যদি $y=x\oplus x$ হয়, তবে $x\oplus n=x\oplus x$; যোগের
  কর্তন-বিধি থেকে $x=n$, যা $x$-কে মানক করে ফেলবে।
  এই বিধি $\Struct{N}$-এ সত্য, তাই $\Struct{M}$-এও।
  $y=x\oplus x\oplus\mathbf{z}^*$ হলেও একই যুক্তি চলে।
\item একই ধরনের যুক্তিতে কোনো বৃহত্তম খণ্ডও নেই।
\item খণ্ডগুলির ক্রম ঘন: $[x]$, $[y]$-এর চেয়ে ছোট হলে
  (যেখানে $x \not\approx y$) তাদের মাঝখানে উভয়ের থেকে
  ভিন্ন একটি খণ্ড~$[z]$ আছে। ধরা যাক $x\prec y$।
  আগের মতো $x\oplus y$-কেও দুই দিয়ে ভাগ করা যায়
  (সম্ভবত ভাগশেষ রেখে)। তাই এমন $u \in \Domain{M}$ আছে,
  যাতে হয় $x\oplus y=u\oplus u$, নয়
  $x\oplus y=u\oplus u\oplus\mathbf{z}^*$।
  উপাদান~$u$ হলো $x$ ও~$y$-এর গড়, তাই তাদের মাঝখানে।
  ধরা যাক $x\oplus y=u\oplus u$; অন্য ক্ষেত্রটি একই রকম।
  $u\approx x$ হলে কোনো মানক~$n$-এর জন্য
  \[
  x \oplus y = x \oplus n \oplus x \oplus n,
  \]
  ফলে $y=x\oplus n\oplus n$ এবং $x\approx y$ হতো,
  যা অনুমানের বিরোধী। অতএব $u\not\approx x$।
  একই ধরনের যুক্তিতে $u\not\approx y$-ও।
\end{enumerate}
অতএব অমানক খণ্ডগুলি মূলদ সংখ্যার মতো ক্রমিত: তারা অন্তবিন্দুহীন
একটি !!a{enumerable} রৈখিক ক্রম গঠন করে। ফলে সত্য পাটিগণিতের যেকোনো দুটি
!!{enumerable} অমানক মডেল $\Struct{M}_1$ ও~$\Struct{M_2}$-এর কেবল
$<$ ও~$=$-যুক্ত ভাষায় রিডাক্টগুলি পরস্পর সমরূপ। বস্তুত
সমরূপতা~$h$ এভাবে সংজ্ঞায়িত করা যায়: $\Struct{M_1}$ ও
$\Struct{M_2}$-এর মানক অংশগুলি মানক মডেল~$\Struct{N}$-এর
সমরূপ, তাই পরস্পরেরও সমরূপ। অমানক অংশ গঠনকারী খণ্ডগুলি
নিজেরাই মূলদ সংখ্যার মতো ক্রমিত; তাই
\olref[dlo]{thm:cantorQ} অনুসারে সেগুলিও পরস্পর সমরূপ।
প্রতিটি খণ্ডে নির্বিচারে উপাদান মিলিয়ে খণ্ডগুলির সমরূপতাকে
খণ্ডের \emph{ভিতরে} বিস্তৃত করা যায়। তারপর $\Struct{M_1}$-এ
$x$-এর উত্তরসূরির প্রতিচ্ছবি নিই $\Struct{M_2}$-এ
$x$-এর প্রতিচ্ছবির উত্তরসূরি। লক্ষ করো, পূর্ণ পাটিগণিতের
ভাষায় $\mathfrak{M}_1$ ও~$\mathfrak{M}_2$ সমরূপ—এ কথা
\emph{অনুসৃত হয় না}; সমরূপতা সব সময়ই কোনো সংকেতসমষ্টির সাপেক্ষে।
কারণ $\Domain{M_1}$ ও~$\Domain{M_2}$-এর উপর যোগ ও গুণ
সংজ্ঞায়িত করার অসমরূপ উপায় আছে। (ভটের একটি বিখ্যাত উপপাদ্য
থেকেও এটি অনুসৃত হয়: কোনো সম্পূর্ণ তত্ত্বের গণনীয় মডেলের
সংখ্যা~2 হতে পারে না।)

\begin{prob}
দেখাও, পাটিগণিতের কোনো অমানক মডেলে বৃহত্তম খণ্ড থাকতে পারে না।
\end{prob}

\begin{prob}
$\Lang{L}$ প্রথম-ক্রমের এমন !!{language} হোক, যাতে
$\eq$ ছাড়া একমাত্র !!{predicate}~$<$; আর
$\Struct{N} = (\Nat, <)$ হোক। $\Struct{N}$-এর সব
সসীম ও সসীম-পরিপূরক উপসেট সংজ্ঞায়নযোগ্য। দেখাও, কেবল
এই উপসেটগুলিই $\Struct{N}$-এ সংজ্ঞায়নযোগ্য।

(ইঙ্গিত: প্রথমে $\Obj{prc}(x,y)$ দিয়ে নিচের
$\Lang{L}$-সূত্রটি বোঝাই, যার অর্থ “$x$ হলো $y$-এর
অব্যবহিত পূর্বসূরি”:
\[
x<y \land \lnot \lexists[z][(x<z \land z < y)].
\]
এবার $\Struct{N}$-এর প্রতিটি সংজ্ঞায়নযোগ্য উপসেটের জন্য
$\Lang{L}$-এর একটি সূত্র~$!A(x)$ আছে। এমন যেকোনো $!A$-এর
জন্য নিচের বাক্য~$!D$ বিবেচনা করো:
\[
\lexists[x][\lforall[y][\lforall[z][((x<y \land x<z \land \Obj{prc}(y,z)
  \land !A(y)) \lif !A(z))]]].
\]
দেখাও, $\Sat{N}{!D}$ যদি এবং কেবল যদি $!A$ দ্বারা
সংজ্ঞায়িত $\Struct{N}$-এর উপসেটটি সসীম অথবা তার পরিপূরক
সসীম।

এবার $\Struct{N}$-এর সঙ্গে প্রাথমিক সমতুল্য একটি অমানক
মডেল~$\Struct{M}$ নিই। $a \in \Domain{M}$ অমানক হলে
$a$-এর চেয়ে বড় $b,c \in \Domain{M}$ নিই, যেখানে
$b$ হলো~$c$-এর অব্যবহিত পূর্বসূরি। তাহলে
$\Domain{M}$-এর এমন একটি স্বসমরূপতা~$h$ আছে যাতে
$h(b)=c$ (কেন?)। ফলে $b$ যদি $!A$ পরিতৃপ্ত করে, তবে
$c$-ও করে (কেন?)। তাই $!D$, $\Struct{M}$-এ সত্য;
সুতরাং $\Struct{N}$-এও সত্য। কিন্তু এর অর্থ,
$!A$ দ্বারা সংজ্ঞায়িত $\Struct{N}$-এর উপসেটটি সসীম
অথবা সসীম-পরিপূরক।)
\end{prob}

\end{document}
